\begin{exercisechunk}
\item \textbf{Finite precision produces an eventually periodic state.}\suppmark
\label[exercise]{ex:l14-finite-precision}
\exerciseprofile{2}{1}{2}{\exproof\quad\exinterpretation}{%
Identify the finite-state limitation created by fixed memory and fixed
precision.}
Turing-completeness claims depend on the available memory or numerical
precision. This exercise proves a limitation that follows from fixed finite
precision without using any architecture-specific lower bound.

Consider a deterministic recurrent system with \(d\) state coordinates, each
stored with exactly \(b\) bits. The system reads the unary input
\(a^t=a\cdots a\) one symbol at a time and produces one output bit after each
update. The state-update rule and output rule are fixed for all input lengths.
\begin{enumerate}[(a),beginpenalty=10000]
\item Prove that the recurrent system has at most \(2^{bd}\) internal states.
Deduce that its state sequence on the infinite unary stream
\(a,a,a,\ldots\) is eventually periodic, with both the preperiod and period
bounded by \(2^{bd}\).
\item Deduce that the output sequence is eventually periodic.
\item Let
\[
 \cL_{\mathrm{pow2}}=\{a^{2^k}:k\ge0\}.
\]
Prove that the indicator sequence
\(y_t=\I(a^t\in\cL_{\mathrm{pow2}})\) is not eventually periodic. Conclude
that no recurrent system satisfying the stated fixed-precision assumptions
recognizes \(\cL_{\mathrm{pow2}}\) correctly at every input length.
\item Explain which step of the proof fails when the system has an external
unbounded tape or a state coordinate whose available precision can grow with
the execution length.
\end{enumerate}
\end{exercisechunk}
